\chapter{اللقاء الحادي عشر: تابع تفسير سورة البقرة وتأصيل منهجي}

\section{مراجعة لمنهج الطبري مع الإسرائيليات}
السلام عليكم ورحمة الله وبركاته. الحمد لله رب العالمين، والصلاة والسلام على نبينا محمد وعلى آله وصحبه أجمعين.

نبدأ بمراجعة لخلاصة منهج \rawi{الطبري} في التعامل مع الإسرائيليات كما مر معنا في الآيات السابقة (مثل قصة زلة آدم). \rawi{الطبري} يورد الإسرائيليات لبيان ما أجمل في الآيات، ولكنه قد يقبل بعضها إذا وافق ما في كتاب الله، ويرد ما خالف كتاب الله أو خالف لغة العرب أو خالف إجماع المفسرين. ويتوقف فيما لا يصدق ولا يكذب.

\begin{fawaid}[قاعدة في إطلاق المنهج]
حينما تتكلم عن "منهج" عالم في مسألة، فهذه كلمة كبيرة تدل على أنك استقرأت كل ما في كتابه ووصلت إلى نتيجة عامة. والأصوب لمن يدرس موضعاً معيناً أن يقول: "الذي يظهر من هذا الموضع كذا وكذا"، ولا يطلق القول بالمنهج العام إلا بعد الاستقراء التام.
\end{fawaid}

\separator

\section{تأويل (يا بني إسرائيل اذكروا نعمتي التي أنعمت عليكم)}
\isnad{قال الإمام الطبري رحمه الله:} «القول في تأويل قوله: \begin{ayah}يَا بَنِي إِسْرَائِيلَ\end{ayah}. إسرائيل بمعنى عبد الله، وصفوته من خلقه... وانما خاطب الله جل وعز بقوله يا بني إسرائيل أحبار اليهود من بني إسرائيل الذين كانوا بين ظهراني مهاجر رسول الله صلى الله عليه وسلم...».

وبين \rawi{الطبري} أن سبب الخطاب لهم وذكر أنباء أسلافهم التي لا يعلمها إلا هم، هو ليكون حجة للنبي صلى الله عليه وسلم وآية على صدقه؛ لأنه لم يصل إلى ذلك إلا بوحي من الله. 

\isnad{وأما تأويل:} \begin{ayah}اذْكُرُوا نِعْمَتِيَ الَّتِي أَنْعَمْتُ عَلَيْكُمْ\end{ayah} فهو اصطفاؤه منهم الرسل وإنزال الكتب عليهم واستنقاذهم من آل فرعون. والله يذكر نعمه لتكون دليلاً على ربوبيته وقدرته، ولتكون حجة ليُعبد ويُشكر.

\separator

\section{تأويل (وأوفوا بعهدي أوف بعهدكم) واستدعاء النظائر}
\isnad{قال أبو جعفر:} «عهد الله ووصيته التي أخذ على بني إسرائيل في التوراة أن يبينوا للناس أمر محمد صلى الله عليه وسلم أنه رسول الله... وأوفي بعهدكم: أدخلكم الجنة».

\begin{fawaid}[استدعاء النظائر وتفسير القرآن بالقرآن]
من منهج الطبري أنه يذكر "النظائر" القرآنية لتفسير الآية. والنظير يُستدعى لغرضين:
\begin{itemize}
    \item \textbf{إما لتفسير الآية بالآية:} كاستدعائه لقوله تعالى: \textit{ولقد أخذ الله ميثاق بني إسرائيل وبعثنا منهم اثني عشر نقيبا} لبيان ما هو العهد المذكور في هذه الآية.
    \item \textbf{وإما لبيان ما تندرج تحته الآية من عموم:} كاستدعائه لقوله: \textit{إن الله اشترى من المؤمنين أنفسهم... ومن أوفى بعهده من الله} لبيان قاعدة الوفاء بالعهد عموماً.
\end{itemize}
\end{fawaid}

\isnad{ثم قال الطبري في تأويل:} \begin{ayah}وَإِيَّايَ فَارْهَبُونِ\end{ayah}: «إياي فاخشوا واتقوا أيها المضيعون عهدي...».

\begin{fawaid}[الفرق بين النزول والتنزل]
يفرق المحققون بين "نزول" الآية و"تنزلها"؛ فهذه الآية نزلت في أحبار اليهود، لكنها تتنزل على كل من كان عنده علم وكتمه، وتشمل كل مؤمن بأن يرهب ربه. فالصواب أن نبدأ من النزول التاريخي للآية ولكن لا نقف عنده.
\end{fawaid}

\separator

\section{تأويل (وآمنوا بما أنزلت) والفرق بين صحة المعنى وصحة التفسير}
\isnad{قال أبو جعفر:} «يعني بقوله \textit{وآمنوا}: صدقوا... ويعني بقوله \textit{بما أنزلت}: ما أنزل على محمد من القرآن. \textit{مصدقا لما معكم}: لما مع اليهود من التوراة».

\subsection{تأويل (ولا تكونوا أول كافر به)}
اختلفوا في الضمير (به) على ماذا يعود؟ 
قال \rawi{الطبري}: «ولا تكونوا أول أمتكم كذب به (أي بالقرآن)». وذكر قراءات أخرى أن الضمير يعود على (محمد صلى الله عليه وسلم) أو على (التوراة). ورجح الطبري أنه يعود على القرآن لأن أول الآية أمرت بالإيمان بما أُنزل (وهو القرآن) فكان الواجب أن يكون النهي في آخرها عن الكفر به.

\begin{fawaid}[صحة المعنى لا تستلزم صحة التفسير]
من قال إن اليهود بكفرهم بالقرآن قد كفروا بمحمد وبكتابهم، فكلامه ومعناه صحيح في ذاته بلا شك. ولكن الطبري يرد هذا القول كـ "تفسير" لهذه الآية المعينة، لأن السياق يرجح عودة الضمير على القرآن، فصحة المعنى العام لا تعني بالضرورة أنه هو المراد من الآية.
\end{fawaid}

\separator

\section{تأويل (ولا تشتروا بآياتي ثمنا قليلا وإياي فاتقون)}
\isnad{قال أبو جعفر:} «لا تبيعوا ما آتيتكم من العلم بكتابي وبآياتي بثمن خسيس وعرض من الدنيا قليل... وهو رضاهم بالرياسة على أتباعهم... \textit{وإياي فاتقون}: يقول فاتقوني في بيعكم آياتي بالخسيس من الثمن، وكفركم بما أنزلت».

\begin{fawaid}[تخصيص دلالة الألفاظ العامة بالسياق]
من بديع منهج \rawi{الطبري} أنه لا يبقي الألفاظ العامة (مثل: التقوى، الظلم، البر) على عمومها المطلق حين يفسر الآية، بل يعين معناها بما يناسب السياق. فبدلاً من أن يفسر (فاتقون) باجتناب المنهيات عموماً، فسرها هنا بـ: "فاتقوني في بيعكم آياتي"، وهذا هو التفسير الدقيق الذي يربط اللفظ بمحله.
\end{fawaid}

\separator

\section{تأويل (ولا تلبسوا الحق بالباطل وتكتموا الحق وأنتم تعلمون)}
\isnad{قال أبو جعفر:} «ولا تلبسوا: لا تخلطوا... فكان لبس المنافق منهم الحق بالباطل: إظهاره الحق بلسانه إقراره بمحمد... وخلطه ذلك الظاهر بالباطل الذي يستبطنه. وكان المقر منهم بأنه مبعوث إلى غيرهم الجاحد أنه مبعوث إليهم، خلطاً للحق بالباطل».

ثم بين الطبري وجوهاً إعرابية لقوله: \begin{ayah}وَتَكْتُمُوا الْحَقَّ\end{ayah} إما مجزومة عطفاً على (تلبسوا)، وإما منصوبة (على اعتبار الواو للمعية)، أي: لا تلبسوا الحق بالباطل مع كتمانكم الحق.

\separator

\section{تأويل (وأقيموا الصلاة وآتوا الزكاة واركعوا مع الراكعين)}
\isnad{قال أبو جعفر:} «وأما الصلاة في كلام العرب فإنها الدعاء... وأرى أن الصلاة المفروضة سميت صلاة لأن المصلي متعرض لاستنجاح طلبته من ثواب الله بعمله، مع ما يسأل ربه فيها من حاجاته».

\begin{fawaid}[استدراك على حصر الصلاة لغوياً في الدعاء]
تأثر بعض اللغويين بفكرة أن الكلمة تُوضع لمعنى ثم تُستعار لغيره، فادعوا أن أصل الصلاة في لغة العرب هو الدعاء المجرد فقط، ثم استُعملت في الشرع للعبادة المخصوصة. والصواب أن الصلاة في لسان العرب أعم من مجرد الدعاء، بل هي معروفة عندهم كعبادة ذات هيئة وتذلل، كما قال تعالى: \textit{وما كان صلاتهم عند البيت إلا مكاء وتصدية}.
\end{fawaid}

\isnad{وأما الزكاة:} فذكر الطبري لها وجهين في كلام العرب: إما من "النماء والزيادة"، وإما من "التطهير". ورجح في زكاة المال أنها بمعنى "التطهير" (تطهير لما بقي من ماله وتخليص له)، وإن كان الوجه الأول صحيحاً أيضاً.

\separator

\section{تأويل (أتأمرون الناس بالبر وتنسون أنفسكم)}
\isnad{قال أبو جعفر:} «وجميع الذي قال في الآية من ذكرنا قوله متقارب المعنى؛ لأنهم وإن اختلفوا في (البر) الذي كان القوم يأمرون به غيرهم... فالتأويل الذي يدل على صحته ظاهر التلاوة: أتأمرون الناس بطاعة الله وتتركون أنفسكم؟».

وهنا أبقى الطبري لفظ (البر) على عمومه ولم يخصصه بصورة واحدة مما ذكره المفسرون، لأنه لا حجة على التخصيص.

\isnad{وأما قوله (وتنسون أنفسكم):} «معنى نسيانهم أنفسهم في هذا الموضع، غير النسيان الذي قال جل ثناؤه: \textit{نسوا الله فنسيهم}، هو بمعنى تركوا طاعته فتركهم الله من ثوابها». فالنسيان هنا بمعنى الترك لا بمعنى غفلة الذاكرة.

\separator

\section{تأويل (واستعينوا بالصبر والصلاة)}
\isnad{قال أبو جعفر:} «تأويل ذلك عندنا أن الله أمرهم بالصبر على كل ما كرهته نفوسهم من طاعة الله... وأصل الصبر منع النفس محابها وكفها عن هواها. ولذلك قيل للمحبوس للقتل: قُتِلَ صبراً (أي حُبس ليُقتل)».

\isnad{فإن قال قائل:} ما معنى الأمر بالاستعانة بالصلاة؟ قيل: إن فيها تلاوة كتاب الله، الداعية آياته إلى رفض الدنيا ونعيمها، المسلية للنفوس، المذكرة بالآخرة. فهي معونة لأهل الطاعة على الجد فيها.

\separator

\section{تأويل (الذين يظنون أنهم ملاقو ربهم)}
\isnad{طرح الطبري إشكالاً:} «كيف أخبر الله عمن وصفه بالخشوع له أنه (يظن) أنه ملاقيه؟ والظن شك؟ قيل: إن العرب قد تسمي اليقين ظناً والشك ظناً، نظير تسميتهم الظلمة سدفة والضياء سدفة».

\begin{fawaid}[قاعدة: السياق يعين دلالات الألفاظ]
الظن يُستعمل في لغة العرب وفي القرآن بمعنى الشك وبمعنى اليقين. ولا يمكن أن يوجد سياق عربي فيه لفظ يحتمل دلالتين إلا ويقوم معه ما يرجح ويحدد المعنى المراد. ففي قوله: \textit{إني ظننت أني ملاق حسابيه} الظن هنا بمعنى اليقين التام.
\end{fawaid}

\subsection{مناقشة نحوية في إعراب (ملاقو ربهم)}
أورد \rawi{الطبري} اختلاف نحاة البصرة والكوفة في سبب سقوط النون من (ملاقون) وإضافتها إلى (ربهم). 
\begin{fawaid}[تحرير موضع الإجماع والاختلاف]
من الدقة المنهجية أن الطبري قبل أن يذكر الخلاف، حرر موضع الإجماع فقال: «لا تدافع بين جميع أهل المعرفة بلغة العرب في إجازة إضافة الاسم... وإسقاط النون»، ثم حرر موضع الخلاف فقال: «وإنما اختلفوا في السبب الذي من أجله أُضيف وأسقطت النون». وهذه طريقة المحققين في عرض المسائل.
\end{fawaid}

\separator

\section{تأويل (يوما لا تجزي نفس عن نفس شيئا ولا يقبل منها شفاعة)}
\isnad{قال أبو جعفر:} «واتقوا يوما لا تقضي نفس عن نفس شيئا، تغني عنها غنى. لأن قضاء الحقوق هنالك من الحسنات والسيئات».

\subsection{الرد على المعتزلة في نفي الشفاعة}
وقد استدل المعتزلة بقوله تعالى: \begin{ayah}وَلَا يُقْبَلُ مِنْهَا شَفَاعَةٌ\end{ayah} على نفي الشفاعة لمرتكب الكبيرة من المسلمين.
\isnad{فرد الطبري عليهم قائلاً:} «هذه الآيات وإن كان مخرجها عاماً في التلاوة، أراد بها خاصاً في التأويل لتظاهر الأخبار عن رسول الله ﷺ أنه قال: \textit{شفاعتي لأهل الكبائر من أمتي}... وأن \textit{ولا يقبل منها شفاعة} إنما هي لمن مات على غير توبة من الكفار. وليس هذا من مواضع الإطالة في القول في الشفاعة، فسنستقصي الحجاج في ذلك في مواضعه إن شاء الله».

\begin{fawaid}[منهج الطبري في تأخير التفصيل لمظانه]
من منهج الطبري أنه يكتفي بذكر الرد الموجز على الفرق المخالفة إذا عرضت الآية، ثم يحيل التفصيل والاستقصاء إلى الموضع القرآني الأليق بالمسألة تجنباً للاستطراد.
\end{fawaid}

\separator

\section{تأويل (وإذ نجيناكم من آل فرعون يسومونكم سوء العذاب)}
بيّن \rawi{الطبري} أن هذا تذكير بنعمة الله عليهم، ثم ساق الإسرائيليات في تفصيل عذاب فرعون وكيف كان يذبح الأبناء ويستحيي النساء.

\subsection{الرد على تأويل (يستحيون نساءكم) بالاسترقاق}
انتقد \rawi{الطبري} قول \rawi{ابن جريج} الذي فرّ من تفسير "الاستحياء" بإبقاء الصبايا أحياء لأن الطفلة لا تسمى (امرأة أو نساء)، ففسرها بـ "يسترقون نساءكم".

\isnad{فرد الطبري قائلاً:} «حاد ابن جريج... فدخل فيما هو أعظم مما أنكر. فتأويله \textit{ويستحيون} بـ \textit{يسترقون} تأويل غير موجود في لغة عربية ولا أعجمية؛ وذلك أن الاستحياء إنما هو استفعال من الحياة نظير الاستبقاء».

\separator

\section{تأويل (وإذ فرقنا بكم البحر) ونقد تكلف اللغويين}
\isnad{قال الطبري:} «ومعنى قوله \textit{فرقنا بكم}: فصلنا بكم البحر... وقد زعم بعضهم (يقصد \rawi{الأخفش}) أن معنى قوله: فرقنا من الماء وبينكم... يريد فصلنا بينكم وبينه وحجزنا حين مررتم فيه. وذلك خلاف ما في ظاهر القرآن... وتفريق البحر بهم على ما وصفنا من افتراق سبله بهم».

وهنا يرد الطبري التكلف النحوي الذي يلجأ إلى تقدير المحذوفات وتضمين الحروف ما لم يدل عليه الظاهر.

\separator

\section{تأويل (وإذ واعدنا موسى أربعين ليلة) والترجيح بين القراءات}
اختلف القراء فقرأ بعضهم (وَاعَدْنَا) وقرأ آخرون (وَعَدْنَا).

\begin{fawaid}[الترجيح بين القراءات بالمعنى]
رجح \rawi{الطبري} أن القراءتين صواب وأنهما متفقتان في المعنى رغم اختلاف اللفظ، لأن المواعدة تقتضي الوعد من الطرفين (فالله وعده المناجاة، وموسى وعده اللقاء والرضا بذلك). وردّ على من أنكر قراءة (وَاعَدْنَا) بحجة أن المواعدة لا تكون إلا بين البشر، مبيناً أن هذا تعنت لا دليل عليه.
\end{fawaid}

\separator

\section{قصة بقرة بني إسرائيل والقاعدة الذهبية في المبهمات}
أطال \rawi{الطبري} في ذكر الروايات الإسرائيلية المتعلقة بقصة ذبح البقرة، وعندما وصل إلى قوله تعالى: \begin{ayah}فَقُلْنَا اضْرِبُوهُ بِبَعْضِهَا\end{ayah} واختلافهم في أي جزء من البقرة ضُرب به القتيل (فخِذها أم عظم منها أم غيره).

\begin{fawaid}[منهج الطبري في تفاصيل الإسرائيليات والمبهمات]
قرر \rawi{الطبري} قاعدته الجليلة قائلاً: «الصواب... أن يُقال أمرهم الله أن يضربوا القتيل ببعض البقرة ليحيا... ولا خبر تقوم به حجة على أي أبعاضها التي أُمِروا أن يضربوه به... \textbf{ولا يضر الجهل بأي ذلك ضُرِب القتيل، ولا ينفع العلم به} مع الإقرار بأن القوم قد ضربوا القتيل ببعض البقرة فاحياه الله». وهذه قاعدة ذهبية في التعامل مع مبهمات القصص القرآني التي لم يفصلها الوحي؛ فالعبرة تتحقق بالقدر المذكور، وما سواه من الغيب الذي لا طريق للعلم اليقيني به.
\end{fawaid}

\separator

\section{تأويل (فبدل الذين ظلموا قولا غير الذي قيل لهم)}
تكلم \rawi{الطبري} عن تبديل بني إسرائيل لقول (حطة) إلى (حنطة)، ثم انتقل لقوله تعالى: \begin{ayah}فَأَنزَلْنَا عَلَى الَّذِينَ ظَلَمُوا رِجْزًا مِّنَ السَّمَاءِ\end{ayah}.

وساق اختلافهم في معنى (الرِّجْز)، فقيل هو العذاب، وقيل هو الطاعون، مستدلين بحديث: \begin{hadith}إن الطاعون رجز أنزل على من كان قبلكم\end{hadith}.

\begin{fawaid}[عدم القطع في موضع الظن]
مع أن الطبري مال إلى أن الرجز هنا هو الطاعون استئناساً بالحديث، إلا أنه قال: «وإن كنت لا أقول إن ذلك كان يقيناً، لأن رسول الله ﷺ لا بيان فيه أي أمة عُذبت بذلك». وهذا من إنصافه ودقته؛ فاحتمال أن يكون الطاعون نزل على أمة أخرى غير هؤلاء يمنع من القطع والجزم بتفسير الآية به.
\end{fawaid}

\separator

بهذا نكون قد أتممنا بحمد الله المجلد الأول من تفسير الإمام الطبري (وفق طبعة هجر/علم الكتب)، وسنبدأ في الدرس القادم بإذن الله في المجلد الثاني. والسلام عليكم ورحمة الله وبركاته.